% Tilstrekkelig beskrivelse til at forsøk og undersøkelser kan reproduseres og etterprøves:
% - Forsøksbetingelser, og hvordan disse stemmer med prinsipper beskrevet i kapittel 2.
% - Praktisk fremgangsmåte.
% - Analysemetode (oftest i vedlegg).

\chapter{PERFORMANCE TESTING}

Since Snowman is similar to BHTB and Snortmanager, it has been important to
benchmark their performance in various formal test cases in order to
highlight the strength and weaknesses in each software solution. The
purpose of the tests is to reach final conclusions on key questions regarding
the performance of Snowman, and to point out where the difference between the
three software packages lies.

To get a picture of how the software performs on different systems, it was
tested on three architectures (table \ref{tab:testSystems}). To minimize error, each test
was run five times where the mean value of the runs is considered the valid test
result. All tests can be reproduced with the test procedure documentation in
appendix. %TODO: Appendix.

\section{Definitions}

The testing is split up into \textit{test suites}, which contain a
\textit{problem} in the form of a question. The problem forms a clear
description of what is to be tested, and how results should be evaluated.
A test suite also contains one or more test cases which exist to address
the problem. Test cases describe 

each test suite consists of one or 

results from test cases are referred to as the test results.

Each test case is composed of a series of 



\section{Environment}
%Beskriv testoppsett. maskinvare, nettverk, programvare, etc.
All tests are carried out on three systems, named \textit{S1}, \textit{S2} and
\textit{S3} (table \ref{tab:testSystems}). These systems differ in processing
power, available memory and storage technology, and the results will indicate
what impact these factors have for software performance. The performance
testing in this report is also prone to error, as tests might get disturbed by
other background processes etc. Testing on three systems gives results with a
higher confidence, as results are checked aginst each other for 
inconsistencies. When an inconsistency in results is found, the test is run
again to see if the inconsistency prevails.

%Set row spacing:
\renewcommand{\arraystretch}{1.5}

\begin{table}[h!]
  \centering
  \begin{tabular}{l|l|l|l}
  	%& \multicolumn{1}{c|}{Test} &
  	%\multicolumn{1}{c|}{S1} & \multicolumn{1}{c|}{S2} & \multicolumn{1}{c}{S3}
  	% \\
	System & Processor type & Memory & Storage \\
	\hline
	S1 & Intel\textregistered{} Core\texttrademark{} 2 Solo, 1.4GHz & 4GB RAM & 5K
	RPM HDD \\
	\hline
	S2 & Intel\textregistered{} Core\texttrademark{} i7, 6x3.2GHz & 12GB RAM & 10K
	RPM HDD \\
	\hline
	S3 & Intel\textregistered{} Core\texttrademark{} i5, 4x3.2GHz & 16GB RAM & SSD
	\\
	\hline
	\multicolumn{4}{l}{\textit{Note: See appendix for full details.}}
  \end{tabular}
  \caption{Test systems.}
  \label{tab:testSystems}
\end{table}

Deviation

\section{Test suites}

\subsection{Update processing}
The Update module is tested for speed and linearity, which is compared to  
\subsubsection{Test 1}
\paragraph{Problem: How fast can the software process an update?}
The purpose of this question is to determine the average
processing time per rule. Results are ranked from short to long
processing time, where shorter is better.
\paragraph{Test cases}
To address this problem, five test cases were introduced. The tests contained
an artificially generated update with 1, 10, 100, 1000 and 10000 rules,
respectively. Rules were split into files with ten each (table
\ref{tab:test1Names}).

%Set row spacing:
\renewcommand{\arraystretch}{1.5}

\begin{table}[h!]
  \centering
  \begin{tabular}{l|l|l}
  	%& \multicolumn{1}{c|}{Test} &
  	%\multicolumn{1}{c|}{S1} & \multicolumn{1}{c|}{S2} & \multicolumn{1}{c}{S3}
  	% \\
	Test case & Rules & Files \\
	\hline
	1-1-1 & 1 & 1 \\
	10-1-10 & 10 & 1 \\
	100-10-10 & 100 & 10 \\
	1000-100-10 & 1000 & 100 \\
	10000-1000-10 & 10000 & 1000 \\
	\hline
  \end{tabular}
  \caption{Test cases in Test 1.}
  \label{tab:test1Names}
\end{table}

\subsubsection{Test 2}
\paragraph{Q2: How strong is the linear correlation between processing time and
the size of an update?}
This question rises from the requirements specification of Snowman (chapter
\ref{chap:requirements_spec}), where it is stated that Snowman shall exhibit a
linear increase in processing time in correlation to the number of rules in the
update (where the number of rules depicts the update size). Results are ranked
from strong correlation to no correlation, where stronger is better.
% TODO: remeber to explain, either here or in the rec. spec., why this is an
% important measure, e.g. no coorelation might imply programming errors.
 
\paragraph{Q3: Does raising the ratio between files and rules increase
processing time?}
Updates can come with the rules either gathered in one file, or distributed over 
several files. This question is asked to examine the difference in processing
time for these two scenarios. Results are ranked from no difference to big
difference, where less difference is better.

%\paragraph{Q4: Is the processing time identical for active and inactive
%rules?}
%Assess whether active and inactive rules are parsed differently..?
% -do we really need this? We can just check the source code.
% Question abandoned.

\paragraph{Q4: Is the processing time shorter when the update contains no new
rules?}


\paragraph{Q5: Are the processing times for two updates identical when only one
rule has changed between the updates?}
This question seeks to determine how the software handles an update when there
is only one change from the previous update. This is an important factor to
consider, since an update usually contains only a small amount of changes from
previous updates. Results are ranked from short processing time for second
update to no difference in processing time, where shorter time is better.

%TODO: when analyzing results, be sure to highligt that a completely new
% update is run once, and subsequent updates often contain a small amount 
% of changes. 









An automatically generated update is used for Q1 through Q6